\vspace{-0.5em}
\section{Discussion}\label{sec:discussion}
\vspace{-0.2em}
\subsection{Discussion of PaTAS Findings}

Beyond per-output assessment, PaTAS can estimate a static trust opinion of the network as \(T(\text{NN}) = PaTAS_{FF}((1,0,0))\), on the principle that a trustworthy model preserves trust; under this definition the IPTA results of Experiment~\ref{exp:3} show clear degradation for adversarial samples (\cref{tab:mnistpoisipta}).

Although trust mass and accuracy are often correlated, our results reveal meaningful divergences. In the breast cancer task (\cref{tab:cancer_results}), the mixed profile \((0.25,0,0.75)\) yields a higher final trust mass (\(0.35\)) than the uncertain-features/trusted-label case (\(0.32\)), yet the latter achieves better test accuracy (96\% vs.\ 94\%); starkest, in poisoned MNIST (\cref{tab:mnistpoisipta}) clean digits \(3\) and \(6\) reach comparable accuracy (97.8\% and 97.4\%) but IPTA trust \(0.886\) vs.\ \(0.570\). Higher trust mass therefore does not imply superior predictive performance, and equal accuracy does not imply equal trustworthiness.

These divergences reflect that PaTAS evaluates the reliability of the inference path, not predictive accuracy alone: digit \(6\) was a poisoned class during training, so the parameters predicting it are less trusted than those for \(3\), even on clean test samples. 

Interpretation of trust values, like accuracy, is application-dependent: \((0.4,0,0.6)\) may be inadequate in high-stakes settings yet sufficient elsewhere. PaTAS supports such reading with a single interpretable metric that complements accuracy.

\paragraph*{Obtaining input and label trust}\label{sec:acquisition}
Three routes provide the input- and label-level trust opinions consumed by PaTAS, in decreasing order of information. (i) \emph{Provenance-based assessment}: when data are drawn from sources of known reliability, the quantification models of \cref{sec:back} map the supporting and contradicting evidence directly into a binomial opinion; this is how our companion work assesses dataset trustworthiness~\cite{10.1007/978-3-032-19096-3_17} and how feature-level opinions are assigned when part of an input, such as a region covered by an adversarial patch, is known to be unreliable. (ii) \emph{Estimation from data}: label agreement across annotators or models, and input-quality signals such as out-of-distribution or anomaly scores, can be converted to opinions through the same evidence mapping; the conformity-based construction of \cref{exp:4} instantiates this route. (iii) \emph{Vacuous fallback}: when no such information exists, the opinion is initialized to \((0,0,1)\); since a vacuous input yields a vacuous output (\cref{thm:infvac}), this fallback conservatively produces uncertainty-dominated assessments rather than unwarranted trust, and PaTAS still contributes through parameter and path trust. Crucially, the per-query collection circumstances of a single input are usually easier to establish than the provenance of an entire training set.

\paragraph*{Magnitude blindness, scale, and cost}
The Trust Function combines opinions independently of weight and activation magnitudes: trust deliberately tracks the provenance and reliability of the quantities entering a computation rather than their functional influence, so an input the network effectively ignores can still contribute to the output opinion. The softmax-weighted IPTA aggregation reported alongside \cref{exp:3} is a first step toward influence-aware readings, and magnitude-aware trust functions are future work. For convolutional layers, the natural extension keeps one opinion per filter weight, shared across spatial positions, with the gradient evidence of all positions fused per batch (our prototype convolutional implementation follows this design); attention layers remain open. On cost, storing one binomial opinion per parameter takes three scalars per weight, an approximately threefold memory overhead matching the analytical compute estimate; measured wall-clock overhead is reported with the revision experiments.

A key factor is the threshold \(\epsilon\) separating positive from negative gradient evidence: as \cref{tab:cancer_results} shows, too small an \(\epsilon\) makes PaTAS read gradients as large even when the model performs well, so that trusted labels reduce parameter trust, resembling the \(g \to +\infty\) case in Appendix~\ref{app:asymptotic}. \(\epsilon\) can be tuned per layer or tied to the learning rate, decaying alongside it to track the expected gradient scale. An explicit sweep over three orders of magnitude (Appendix~\ref{app:ablations}) shows a broad plateau: every reported metric is unchanged for $\epsilon \in [0.05, 1]$ and only degenerate-small values shift the operating point, so $\epsilon$ behaves as a scale parameter with a wide safe range rather than a tuned hyperparameter. The revision-fusion variant is similarly immaterial (differences at most $0.002$ across average, cumulative, and weighted fusion, Appendix~\ref{app:ablations}); averaging is retained because Theorem~\ref{thm:arithm} guarantees non-degenerate long-run behavior under repeated revision, whereas cumulative fusion drives uncertainty monotonically toward dogmatic saturation, and the independence rationale for cumulative fusion concerns feedforward fusion of feature opinions, a distinct role.

\vspace{-0.5em}
\subsection{Threat Model and Security Scope}\label{sec:threat}
\vspace{-0.5em}

\paragraph*{Assets and goal}
The asset PaTAS protects is the \emph{reliability of individual model predictions at inference time}: the defender wishes to know, per query, whether an output can be trusted given the input quality, the training-data reliability, and the parameters through which the input propagates. PaTAS is a \emph{detection and monitoring} mechanism, not a hardening one; it does not alter the model or prevent an attack, but exposes predictions with weak supporting evidence.

\paragraph*{Adversary}
We consider a \emph{training-time data-poisoning adversary}~\cite{chen2017targetedbackdoorattacksdeep} who can inject or relabel a bounded fraction of training samples, including backdoor triggers such as fixed pixel patches, to induce targeted misclassifications hidden behind high nominal accuracy and confidence. The adversary controls neither the architecture, nor the training procedure, nor the trust assessments of known-provenance data. This is the threat evaluated in \cref{exp:3}.

\paragraph*{Defender knowledge}
PaTAS assumes white-box access to the model (parameters, activations, gradients), as it runs alongside training and inference, and access to trust opinions on inputs and labels; when unavailable, these are initialized to the vacuous opinion \((0,0,1)\), yielding conservative, uncertainty-dominated assessments rather than false assurance (\cref{sec:acquisition}).

\paragraph*{Out of scope}
Evasion attacks crafted at inference time, model-stealing, and membership-inference attacks are outside the present scope; PaTAS's signals may apply to them, e.g., by monitoring the trust profile of query streams, but we make no such claim here. We likewise prove consistency properties of the trust computation (\cref{sec:ptastheory}) rather than formal robustness guarantees against a bounded adversary, which remain an open direction.
